\documentclass[11pt,a4paper]{article}
\usepackage[utf8]{inputenc}
\usepackage[T1]{fontenc}
\usepackage{geometry}
\usepackage{xcolor}
\usepackage{listings}
\usepackage{fancybox}
\usepackage{tcolorbox}
\usepackage{amsmath}
\usepackage{amsfonts}
\usepackage{amssymb}

% Page setup
\geometry{margin=2cm}

% Define colors for syntax highlighting
\definecolor{terminalgreen}{RGB}{0,128,0}
\definecolor{terminalblue}{RGB}{0,0,255}
\definecolor{terminalred}{RGB}{255,0,0}
\definecolor{terminalgray}{RGB}{128,128,128}
\definecolor{backgroundgray}{RGB}{245,245,245}
\definecolor{bordergray}{RGB}{0,0,0}

% Custom style for terminal commands
\lstdefinestyle{terminal}{
    backgroundcolor=\color{backgroundgray},
    basicstyle=\ttfamily\small,
    breakatwhitespace=false,
    breaklines=true,
    captionpos=b,
    commentstyle=\color{terminalgray}\itshape,
    escapeinside={\%*}{*},
    extendedchars=true,
    frame=single,
    frameround=tttt,
    framerule=1pt,
    rulecolor=\color{bordergray},
    keepspaces=true,
    keywordstyle=\color{terminalgreen}\bfseries,
    language=bash,
    morekeywords={sudo, apt, git, cmake, python3, mkdir, cd, source, pip, install},
    numberstyle=\tiny\color{terminalgray},
    showspaces=false,
    showstringspaces=false,
    showtabs=false,
    stepnumber=1,
    stringstyle=\color{terminalred},
    tabsize=2,
    title=\lstname,
    xleftmargin=10pt,
    xrightmargin=10pt,
    aboveskip=15pt,
    belowskip=15pt
}

\begin{document}

\title{\textbf{LLaMA.cpp Installation and Environment Setup}}
\author{Step-by-Step Installation Guide}
\date{\today}

\maketitle

\section{Complete Installation Script}

The following commands will install LLaMA.cpp and set up the required Python environment on Ubuntu/Debian systems:

\begin{lstlisting}[style=terminal]
# Step 1: Update system packages and install build dependencies
sudo apt update && sudo apt install -y build-essential git cmake python3 python3-venv python3-pip curl

# Step 2: Create project directory and navigate to it
mkdir -p ~/llm-local && cd ~/llm-local

# Step 3: Clone the LLaMA.cpp repository and enter the directory
git clone https://github.com/ggerganov/llama.cpp.git && cd llama.cpp

# Step 4: Create build directory and configure the project
mkdir -p build && cd build

# Step 5: Configure CMake with Release build type (CURL disabled for cleaner build)
cmake .. -DCMAKE_BUILD_TYPE=Release -DLLAMA_CURL=OFF

# Step 6: Build the project using all available CPU cores
cmake --build . -j$(nproc)

# Step 7: Return to project root and create Python virtual environment
cd ../.. && python3 -m venv venv && source venv/bin/activate

# Step 8: Upgrade pip and install Python dependencies
python -m pip install --upgrade pip && pip install llama-cpp-python tqdm datasets
\end{lstlisting}

\section{Installation Breakdown}

\subsection{System Dependencies}
\begin{itemize}
    \item \texttt{build-essential}: Essential compilation tools (gcc, g++, make)
    \item \texttt{git}: Version control system for cloning repositories
    \item \texttt{cmake}: Build system generator for C/C++ projects
    \item \texttt{python3}: Python interpreter for running scripts
    \item \texttt{python3-venv}: Virtual environment support
    \item \texttt{python3-pip}: Python package installer
    \item \texttt{curl}: HTTP client for downloading files
\end{itemize}

\subsection{Build Configuration}
\begin{itemize}
    \item \texttt{Release}: Optimized build configuration for production use
    \item \texttt{DLLAMA\_CURL=OFF}: Disables CURL dependency for simpler builds
    \item \texttt{-j\$(nproc)}: Uses all available CPU cores for parallel compilation
\end{itemize}

\subsection{Python Environment}
\begin{itemize}
    \item \texttt{llama-cpp-python}: Python bindings for LLaMA.cpp
    \item \texttt{tqdm}: Progress bars for long-running operations
    \item \texttt{datasets}: Hugging Face datasets library for model data
\end{itemize}

\section{Post-Installation}

After successful installation, the following executables will be available in the \texttt{build} directory:
\begin{itemize}
    \item \texttt{llama-cli}: Command-line interface for text generation
    \item \texttt{llama-server}: HTTP server for API endpoints
    \item \texttt{llama-bench}: Performance benchmarking tool
    \item \texttt{llama-perplexity}: Model quality evaluation tool
\end{itemize}

To activate the Python environment in future sessions, navigate to the project directory and run:
\begin{lstlisting}[style=terminal]
source venv/bin/activate
\end{lstlisting}

\end{document}